\chapter{Background \& Technical Foundations}
\label{ch:background}

This chapter reviews the systems, control mechanisms, and machine learning techniques 
upon which this thesis builds. \Cref{sec:background:scion} introduces the SCION 
path-aware Internet architecture and the path-selection opportunities it creates. 
\Cref{sec:background:qoe} covers quality-of-experience metrics for real-time video, adaptive 
bitrate control, and multipath transport mechanisms. \Cref{sec:background:rl} recalls 
the core reinforcement learning (RL) machinery: Markov decision processes, deep $Q$-networks 
with standard refinements, soft actor-critic, and permutation-equivariant scoring architectures 
for dynamic action spaces. Finally, \Cref{sec:background:intent} details goal-conditioned and
multi-objective policy representations, the machinery that lets one policy serve many objectives.

\section{The SCION Path-Aware Architecture}
\label{sec:background:scion}

SCION~\cite{Perrig2017} is a next-generation, path-aware Internet architecture designed for high
availability and transparency of the forwarding path. Its defining departure from today's Internet
is that \emph{end hosts, not intermediate routers, choose the path a packet takes}. The network is
organized into \emph{isolation domains} (ISDs), each governed by a set of core autonomous systems
(ASes). Path segments are discovered by a beaconing process: core ASes flood \emph{path-segment
construction beacons} (PCBs) that accumulate hop information as they propagate, giving rise to three
segment types -- \emph{up} segments from a leaf AS to its ISD core, \emph{core} segments between
core ASes across ISDs, and \emph{down} segments from a core to a destination leaf. An end host
combines an up, an (optional) core, and a down segment into a complete forwarding path and embeds it
in the packet header.

Two properties of SCION matter for this thesis. First, \textbf{path multiplicity}: a source--
destination pair is typically connected by \emph{several} distinct paths that differ in latency,
available bandwidth, loss, and the ASes they traverse. This is precisely the raw material an
intent-conditioned path-selection agent needs, and the same multiplicity makes multipath transport natural. 
Second, \textbf{path transparency and trust}: because the full AS-level path is visible to the endpoint, 
an agent can reason about properties beyond raw performance -- for instance a \emph{trust} score that
reflects how much confidence to place in a path's advertised or measured quality. We exploit this
in our single-path selection studies, where the reward balances measured goodput against a trust term.

One point of terminology needs fixing before it can mislead. The \emph{trust} term of
\cref{ch:single_path} is a \textbf{performance proxy}, computed from a path's measured loss and delay
alone (\cref{eq:p1trust}). It carries no AS reputation, no certificate, and no path-authentication
content, and it is not trust in the SCION control-plane sense. We keep the name because it is the one
the reward of Bourak~et~al.~\cite{Bourak2025} uses and the one our reward inherits, and we repeat the
qualification where the term is defined.

The flip side of endpoint path choice is a coordination problem: whereas a router runs one routing
process, every SCION host must decide for itself, ideally without flooding the network with probes.
Minimizing that probing overhead is an explicit objective throughout our work.

\subsection{How SCION endpoints select paths today}
\label{sec:background:pathsel}
A SCION host does not receive one path; it receives a set, and something must reduce that set to the
one path a packet carries. In deployed SCION, that reduction is made by \emph{policy} rather than by
measurement. The path service delivers the segments known for a destination, the host composes them
into end-to-end paths, and a path policy -- a declarative filter over hop sequences, ISDs, ASes, and
path length -- selects among the survivors, with ties broken by a static preference such as the fewest
AS hops~\cite{Perrig2017}. Path-aware transport APIs expose the same choice to the application, so that
a socket may pin, rank, or re-select paths. What these mechanisms share is that they rank paths on
\emph{static} properties: neither hop count nor an ISD filter says anything about which candidate is
congested at this hour. That is the status quo \cref{ch:single_path} improves on, and the
\emph{SCION default} baseline of \cref{sec:p1impl:baselines} -- minimum hop count, ties broken by
measured latency -- is our stand-in for it. It is a generous one, and deliberately so: a deployed
default breaks its ties on a static preference, whereas ours is charged a latency probe on every
candidate and breaks them on a measurement.

Two further conventions from that comparison belong here. \emph{Equal-cost multi-path}
(ECMP)~\cite{RFC2992} spreads flows over the equal-cost members of a route by hashing a packet's
five-tuple, which keeps a flow's packets on one member while distributing flows across members; it is
the load-spreading rule the Internet already runs, and \cref{sec:p1impl:baselines} states exactly how
our single-flow evaluation restricts it. And the AS-level topologies the study runs on are generated
with BRITE~\cite{Medina2001} under a Barab\'{a}si--Albert preferential-attachment
model~\cite{Barabasi1999}, which reproduces the heavy-tailed degree distribution of the inter-domain
graph, and then converted into SCION roles and links (\cref{sec:p1design:env}).

\section{Quality of Experience and Transport Mechanisms}
\label{sec:background:qoe}

\subsection{Quality of experience for real-time video}
To evaluate end-to-end intent control, we target interactive video, whose value to a user is captured by 
quality-of-experience (QoE) rather than any single network metric. Perceptual quality at a given encoding 
is commonly summarized by VMAF~\cite{VMAF2016}, a learned metric that predicts subjective video quality; 
interactivity adds extreme sensitivity to end-to-end latency, jitter, and frame loss, since a late frame in a live 
call is effectively a lost frame. Our multipath reward composes a VMAF-based quality term with penalties for 
latency, jitter, and loss, making the QoE objective explicit and tunable.

\paragraph{The ITU-T quality ladder.} The weights in such a reward are the part a reader is entitled to
question, so we anchor ours on published objectives rather than on taste. Four recommendations carry
that provenance and recur throughout \cref{ch:multipath}. G.1010~\cite{ITUG1010} classifies end-user
multimedia by delay tolerance, separating conversational video from one-way video and from
background streaming, which is the axis along which our three application intents are defined.
G.114~\cite{ITUG114} fixes the delay anchors those classes inherit -- one-way delay up to
\SI{150}{\milli\second} is preferred and beyond \SI{400}{\milli\second} is unacceptable for
interactive use. Y.1541~\cite{ITUY1541} supplies the corresponding IP delay-variation objective for its
most stringent class, which is what our jitter normalizer is scaled to. And G.1070~\cite{ITUG1070} is
an opinion model that maps coding bitrate, frame rate, packet loss, and delay to an estimated mean
opinion score for video telephony; because it is an explicit function rather than a table, it can be
used as a reference against which a simpler linear reward is \emph{fit}, which is how the weights of
\cref{tab:p2profiles} are obtained.

\paragraph{Google Congestion Control and the WebRTC stack.} The bitrate half of a deployed real-time
video stack is, in practice, Google Congestion Control (GCC)~\cite{Carlucci2016}, the rate controller
shipped with WebRTC. GCC estimates the available rate from two signals at once: a receiver-side
delay-gradient filter that tracks the one-way delay trend across successive frame groups and reduces the
rate as soon as queues start to build, and a sender-side loss-based controller that reacts to reported
packet loss. The delay-based branch is what makes it a strong opponent for a learned controller in a
real-time setting: it backs off before loss occurs rather than after, which is exactly the behavior a
frame deadline rewards. GCC drives the \texttt{webrtc} comparison point in \cref{ch:multipath} and the
\texttt{path-only} ablation, and is the reference against which the learned rate controller is measured.

\subsection{Adaptive bitrate and learned control}
The application-side lever in joint control is adaptive bitrate (ABR) control: choosing the video encoder's
target rate to match available capacity. Learned ABR controllers such as Pensieve~\cite{Mao2017}
showed that RL can outperform hand-crafted heuristics for buffered video streaming. Our setting differs
in two critical ways: it is strictly real-time (lacking a large playback buffer to hide network variance), 
and the bitrate decision is tightly coupled to a simultaneous \emph{network-side} decision regarding how to 
allocate traffic across multiple paths.

\subsection{Multipath QUIC}
QUIC~\cite{RFC9000} is a modern, encrypted, multiplexed transport running over UDP. Multipath
QUIC~\cite{DeConinck2017} extends it to send a single connection's data across several network paths at
once, in the spirit of Multipath TCP~\cite{RFC8684} but with user-space deployability. Over a 
path-aware substrate like SCION, multiple paths are readily available from path discovery, making MPQUIC 
a natural transport for exploiting SCION's path multiplicity. The open question MPQUIC raises -- 
how to \emph{schedule} bytes across heterogeneous paths -- is one of the primary control dimensions learned 
in this thesis.

\section{Deep Reinforcement Learning Machinery}
\label{sec:background:rl}

\subsection{Markov decision processes}
We model path and rate control as sequential decision-making in a Markov decision process (MDP)
$(\mathcal{S}, \mathcal{A}, P, r, \gamma)$, where at each step the agent observes a state
$s \in \mathcal{S}$, takes an action $a \in \mathcal{A}$, receives a scalar reward
$r(s,a)$, and transitions to $s'$ with probability $P(s' \mid s, a)$. The agent seeks a policy
$\pi$ that maximizes the expected discounted return $\mathbb{E}\!\left[\sum_t \gamma^t r_t\right]$
with discount $\gamma \in [0,1)$. The reward function is where application intent enters: by
changing what the agent is rewarded for, we change what ``best'' means.

\subsection{Deep $Q$-networks and refinements}
The action-value function $Q^\pi(s,a) = \mathbb{E}[\sum_t \gamma^t r_t \mid s_0=s, a_0=a]$ gives the
expected return of taking $a$ in $s$ and following $\pi$ thereafter. A deep $Q$-network
(DQN)~\cite{Mnih2015} approximates the optimal $Q^\star$ with a neural network trained to minimize
the temporal-difference error against a slowly updated \emph{target} network, sampling transitions
from a replay buffer to decorrelate them. We utilize three standard refinements. \emph{Double
DQN}~\cite{VanHasselt2016} decouples action selection from action evaluation in the bootstrap
target, reducing the overestimation bias of vanilla DQN. The \emph{dueling}
architecture~\cite{Wang2016} factors $Q(s,a) = V(s) + \big(A(s,a) - \tfrac{1}{|\mathcal{A}|}\sum_{a'}
A(s,a')\big)$ into a state-value stream $V$ and an advantage stream $A$, which stabilizes learning 
when many paths have comparably good value. \emph{Prioritized experience replay} 
(PER)~\cite{Schaul2016} samples high-error transitions more often, correcting the resulting bias with 
importance-sampling weights.

\subsection{Soft actor-critic}
For continuous control domains -- such as determining target bitrates and continuous traffic split ratios -- 
value-based DQN is ill-suited. We adopt soft actor-critic (SAC)~\cite{Haarnoja2018}, an
off-policy, maximum-entropy actor-critic method. SAC augments the reward with a policy-entropy term,
$\mathbb{E}[\sum_t \gamma^t (r_t + \alpha \mathcal{H}(\pi(\cdot\mid s_t)))]$, which encourages
exploration and robustness; it trains twin $Q$-critics to curb overestimation and automatically tunes the
temperature $\alpha$ against a target entropy. SAC's sample efficiency and stability make it a 
strong fit for expensive, noisy rollouts in network simulators.

\subsection{Variable and structured action spaces}
\label{sec:background:variable}
A central difficulty in network path control is that the number of candidate paths is dynamic: it varies across
source--destination pairs and changes over time. A classical neural network with a fixed output head 
cannot represent this unbounded decision space. The remedy, used throughout this thesis, is a 
\emph{scoring} architecture that is \emph{permutation-equivariant} over the path set: a shared network 
maps each path's feature vector (together with a global context) to a scalar score, and the scores are 
compared -- by $\arg\max$ for discrete single-path selection, or by a softmax to produce a continuous traffic split. 
Invalid or absent paths are masked out. Aggregating a set of per-element embeddings in a way that is
invariant to their ordering is achieved via the Deep Sets construction~\cite{Zaheer2017}, which we employ in 
our continuous critic design. This structural approach decouples the policy from the cardinality of 
the action set, allowing a single network to generalize across topologies with arbitrary numbers of paths.

\section{Goal-Conditioned Reinforcement Learning}
\label{sec:background:intent}

\subsection{Conditioning a policy on its objective}
A policy that must serve diverse application preferences can be made a function of the objective itself. In
goal-conditioned RL, the policy $\pi(a \mid s, g)$ and value function are conditioned on a goal vector $g$,
so that a single neural network represents a whole family of behaviors. Universal value function approximators
(UVFAs)~\cite{Schaul2015} formalize this by learning $V(s, g)$ jointly over states and goals, while hindsight 
experience replay~\cite{Andrychowicz2017} improves sample efficiency by relabeling trajectories with goals 
achieved \emph{a posteriori}. In this thesis, the ``goal'' represents application intent,
encoded as a vector of reward weights: the agent is explicitly told how much it should prioritize throughput
versus latency, jitter, loss, or trust, and is expected to adapt its decisions dynamically.

\subsection{Preference weights as goals: multi-objective RL}
\label{sec:background:morl}
The goal vector used in this thesis is of a particular kind, and naming it correctly places the work in
the right literature. Our intent is not a target state to be reached but a \emph{preference weighting}
over a linear scalarization of several competing reward terms -- goodput against delay in
\cref{ch:single_path}, quality against delay, jitter, and loss in \cref{ch:multipath}. That is the
object of study of \textbf{multi-objective reinforcement learning} (MORL), surveyed by
Roijers~et~al.~\cite{Roijers2013}. Two of its results matter here. First, sweeping the weight simplex
of a linear scalarization traces the convex coverage set of the underlying multi-objective problem, so
``one policy for every intent'' is a statement about representing a family of optima rather than about
storing a list of behaviors. Second, a single network conditioned on the weight vector can represent
that family and generalize to weightings it was not trained on: Yang~et~al.~\cite{Yang2019} learn an
envelope value function over states and weights, and Abels~et~al.~\cite{Abels2019} show a
weight-conditioned network tracking preferences that change during execution -- which is the property
\cref{sec:p1eval:zeroshot} tests for the single-path selector and \cref{sec:p2eval:conditioning} for
the multipath one.

What that literature does not address is \emph{where} in the network the weights must enter, because
its policies emit actions from a fixed, enumerated action space, where any conditioning that reaches
the output layer at all can change the action. Ours are comparison-based, and the difference is
decisive: as \cref{sec:p1design:concat} shows, a conditioning signal can reach a comparison-based
policy, be trained on, and still leave every action unchanged. Hindsight relabeling~\cite{Andrychowicz2017}
is likewise inapplicable to preference goals, since there is no achieved weight vector to relabel with.
The delta this thesis claims on top of MORL is therefore narrow and specific: not that one network can
span a family of objectives, which is established, but where the objective must be injected in a
\emph{scoring} architecture for it to have any effect on the chosen action.

\subsection{Where the conditioning signal enters}
\emph{Where} the goal vector enters the network is as consequential as whether it is supplied at all, and
the literature on goal-conditioned RL says little about it because its usual policies are not
comparison-based. Our policies are. When the action is an $\arg\max$ over per-alternative scores, a
conditioning signal that enters only terms shared by every alternative shifts all scores together and
cancels under the comparison, so the policy can consume the goal without ever acting on it. The signal
must instead reach the per-alternative terms, where it can interact with each candidate's own features
and change their relative order. \Cref{sec:p1design:concat} makes this precise as
\cref{prop:p1cancel}, and \cref{sec:p1eval:ablation} measures its consequences.